\documentclass{article}
\usepackage{enumitem}
\usepackage{amsmath}
\begin{document}

\title{Chapter 21: Neural Control and Coordination}
\date{}
\maketitle

\begin{enumerate}[label=Q\arabic*.]

\item The resting membrane potential of a typical neuron is approximately:
\\(A) +70 mV
\\(B) $-$70 mV
\\(C) 0 mV
\\(D) $-$40 mV

\item During an action potential, depolarization occurs due to:
\\(A) Rapid influx of K$^+$ ions through voltage-gated channels
\\(B) Rapid influx of Na$^+$ ions through voltage-gated Na$^+$ channels
\\(C) Efflux of Na$^+$ ions
\\(D) Influx of Cl$^-$ ions

\item Repolarization during an action potential is caused by:
\\(A) Na$^+$ influx
\\(B) K$^+$ efflux through delayed voltage-gated K$^+$ channels
\\(C) Cl$^-$ influx
\\(D) Na$^+$/K$^+$ ATPase pump activity

\item The refractory period ensures:
\\(A) Neurons can fire at infinite frequency
\\(B) Unidirectional propagation of action potential along the axon
\\(C) Faster conduction velocity
\\(D) Increased neurotransmitter release

\item Assertion: Myelinated nerve fibres conduct impulses faster than unmyelinated fibres.\\
Reason: Saltatory conduction in myelinated fibres allows the action potential to jump between nodes of Ranvier.
\\(A) Both true, Reason is correct explanation
\\(B) Both true, Reason is not correct explanation
\\(C) Assertion true, Reason false
\\(D) Assertion false, Reason true

\item Acetylcholinesterase (AChE) at the synapse functions to:
\\(A) Synthesize acetylcholine in the presynaptic terminal
\\(B) Rapidly hydrolyze ACh after it has acted on the postsynaptic receptor, terminating the signal
\\(C) Transport ACh across the synaptic cleft
\\(D) Store ACh in synaptic vesicles

\item GABA (gamma-aminobutyric acid) is the major:
\\(A) Excitatory neurotransmitter in the brain
\\(B) Inhibitory neurotransmitter in the CNS; it hyperpolarizes the postsynaptic membrane
\\(C) Hormone released from the hypothalamus
\\(D) Neurotrophic factor promoting neuron survival

\item The cerebellum is primarily responsible for:
\\(A) Conscious thought and personality
\\(B) Coordination of voluntary movement, balance, and posture
\\(C) Regulation of hunger and thirst
\\(D) Processing visual information

\item Which part of the brain is the ``master switchboard'' for sensory relay to the cerebral cortex?
\\(A) Hypothalamus
\\(B) Thalamus
\\(C) Medulla oblongata
\\(D) Basal ganglia

\item The knee-jerk reflex (patellar reflex) is a monosynaptic reflex arc involving:
\\(A) Sensory neuron $\rightarrow$ interneuron $\rightarrow$ motor neuron
\\(B) Sensory neuron directly synapsing on motor neuron in the spinal cord
\\(C) Sensory neuron $\rightarrow$ brain $\rightarrow$ motor neuron
\\(D) Two interneurons and one motor neuron

\item Parkinson's disease results from degeneration of dopaminergic neurons in the:
\\(A) Cerebral cortex
\\(B) Substantia nigra (basal ganglia)
\\(C) Cerebellum
\\(D) Hippocampus

\item The blood-brain barrier (BBB) is primarily formed by:
\\(A) Myelin sheaths of axons
\\(B) Tight junctions between brain capillary endothelial cells and astrocyte foot processes
\\(C) The meninges (dura mater)
\\(D) Choroid plexus cells

\item Which lobe of the cerebral cortex processes visual information?
\\(A) Frontal lobe
\\(B) Parietal lobe
\\(C) Temporal lobe
\\(D) Occipital lobe

\item Match the following neurotransmitters with their primary actions:\\
1. Dopamine $\rightarrow$ (a) Reward, movement coordination\\
2. Serotonin $\rightarrow$ (b) Mood regulation, sleep\\
3. Norepinephrine $\rightarrow$ (c) Fight-or-flight arousal\\
4. Glutamate $\rightarrow$ (d) Major excitatory neurotransmitter in CNS
\\(A) 1-a, 2-b, 3-c, 4-d
\\(B) 1-b, 2-a, 3-d, 4-c
\\(C) 1-c, 2-d, 3-a, 4-b
\\(D) 1-d, 2-c, 3-b, 4-a

\item The autonomic nervous system (ANS) controls which of the following?
\\(A) Voluntary skeletal muscle movements
\\(B) Involuntary functions: heart rate, digestion, respiration rate, glandular secretions
\\(C) Conscious sensory perception
\\(D) Language and speech production

\end{enumerate}
\end{document}
